\subsection{Post-solve optimality failures in banded computation}
The banded-Cholesky implementation recorded 9096 method-level
timing observations rejected by the scaled stationarity-residual
check in Eq.~\eqref{eq:kkt}. They represent 3384 observations
under CV, 4160 under GCV, and 1552 under marginal likelihood.
Among the \emph{rejected} records, the smallest reported value
of $\eta_{\mathrm{opt}}$ was $1.0013\times10^{-7}$, the median was
$1.4079\times10^{-7}$, and the largest was
$1.0587\times10^{-6}$, against the specified threshold
$10^{-7}$. Figure~\ref{fig:kktresidualdist} displays this
conditional distribution. The data do not record the
corresponding residual magnitudes for accepted runs; the
histogram must not be used as an estimate of the full
residual distribution or of an independent failure probability.
Since the check follows factorization and solution, the
failures establish rejection by the \emph{complete banded
criterion pipeline}, not an inability of Cholesky to
factor an SPD matrix.

\begin{figure}[t]
\centering
\includegraphics[width=\linewidth]{\KKTDir/kkt_rejection_distribution.pdf}
\caption{Distribution of the scaled first-order optimality
residual $\eta_{\mathrm{opt}}$ among \emph{rejected} banded
complete-criterion computations. Histogram segments are
colored by statistical criterion (CV, GCV and Gaussian
marginal likelihood). The dashed line marks the
implementation-specific $10^{-7}$ acceptance threshold;
the dotted line marks the median of rejected records.
Counts are correlated method-level timing observations,
not independent failures. Residuals of accepted
evaluations are not available in the archived record.}
\label{fig:kktresidualdist}
\end{figure}

A particularly informative contrast occurs in nine parameter configurations. For each configuration, 97 independent realizations were considered. Both methods completed every recorded single-candidate evaluation, but the banded implementation rejected every complete-grid evaluation, whereas the spectral implementation completed the corresponding grids. The configurations comprise CV and GCV with $(n,d,M)=(128,3,1000)$, $(256,4,100)$ and $(256,3,1000)$; GCV with $(64,3,10000)$; and profiled Gaussian marginal likelihood with $(256,3,100)$ and $(512,3,100)$. Figure~\ref{fig:cholgridfail} displays these results.

The nine configurations account for 6984 banded-rejection timing records; another 2112 such records arise from larger fourth-order configurations. These are repeated method-level observations within 97 independent seed clusters per configuration, rather than thousands of independent experimental realizations. The saved status information does not identify the exact smoothing-grid point that first produced each rejection, and hence no candidate-specific instability boundary is inferred.

\begin{figure}[t]
\centering
\includegraphics[width=\linewidth]{\VizDir/figures/fig05_grid_completion.pdf}
\caption{Full-grid numerical completion in nine parameter configurations, with method-specific coloring and patterned rejection fractions. All recorded single-candidate outcomes complete for both implementations, but the banded full-grid pipeline is rejected while the spectral pipeline completes. Each configuration has 97 independent seed clusters; within-seed repeats, mapping policies and multiple-series workloads are correlated.}
\label{fig:cholgridfail}
\end{figure}

\input{\VizDir/tables/table03_grid_failures.tex}

The apparent spectral advantage must be interpreted as \emph{observed computational completion}, not a controlled comparison of identical backward-error tests. The spectral scorer does not impose the same implementation-specific acceptance criterion as the banded scorer. Independent numerical reference calculations would be required to establish a strict relative-accuracy ranking at each penalty.

\subsection{Smoothness inversion remains a common numerical limitation}
An additional failure mechanism affects both implementations at larger difference order and series length. In the fourth-order CV/GCV configurations with $n=1024$, 2048 and 4096, the numerical inversion of the normalized smoothness mapping sometimes returned values outside its prescribed tolerance. Across the evaluated configurations, 6336 timing records explicitly describe inaccurate smoothness inversion: 2112 associated with the banded method and 4224 with SVD. The distribution of observed outcomes in the fourth-order settings is given in Fig.~\ref{fig:fourthfail}.

This error concerns the transformation from the bounded coordinate $S$ to the penalty $\lambda$, rather than evidence that a spectral decomposition failed. Crucially, a successfully completed one-candidate evaluation cannot be used as a substitute for an incomplete multi-candidate search. Accordingly, the large-order configurations with shared inversion difficulties do not supply an empirical winner for the full criterion grid.

\begin{figure}[t]
\centering
\includegraphics[width=\linewidth]{\VizDir/figures/fig06_order4_failures.pdf}
\caption{Stacked method-colored fractions of numerically complete evaluations and distinct rejection mechanisms for difference order four. The bands show implementation-specific complete, post-solve KKT-rejected and inaccurate smoothness-inversion outcomes. Denominators consist of available method-level records; both implementations exhibit inversion failures at higher sample lengths.}
\label{fig:fourthfail}
\end{figure}

\subsection{The single-candidate versus grid-search crossover}
The principal runtime trade-off becomes visible when paired, numerically admissible observations are compared across candidate counts. Figure~\ref{fig:onegrid} shows three examples evaluated under identical matching rules. In conditional CV with $n=256$, $d=3$, the geometric banded/SVD cold-time ratio changes from 0.698 (95\% interval 0.671--0.724) for one candidate to 4.846 (4.567--5.125) for 100 candidates. In fixed-mean GCV with $n=64$, $d=1$, the corresponding ratios are 1.001 (0.982--1.018) and 2.328 (2.282--2.374). For fixed-mean CV with $n=128$, $d=3$, they are 1.069 (1.052--1.087) and 5.401 (5.244--5.555).

\begin{figure}[t]
\centering
\includegraphics[width=\linewidth]{\VizDir/figures/fig07_crossover_distributions.pdf}
\caption{Distributional comparison of matched banded-Cholesky/SVD cold-time ratios for one candidate and 100 candidates. Violin densities, boxplots and individual seed-level points show the pronounced reversal of computational preference. Values above one favor SVD. Corresponding independent-seed bootstrap intervals for geometric means are reported in the text rather than inferred from violin widths.}
\label{fig:onegrid}
\end{figure}

These results support the computational interpretation of Eqs.~\eqref{eq:spectraltrend}--\eqref{eq:rss}: spectral setup is paid once, while the attenuation coefficients and the statistics needed for subsequent criterion evaluations are reused. The crossover is observed in complete scoring workloads; it should not be conflated with the cost of fitting a trend for a prescribed penalty.

\subsection{Memory requirements and computational feasibility}
The principal disadvantage of the spectral approach is its dense storage requirement. For the evaluated implementation, the conservative memory admission allowance corresponds to five $n\times n$ double-precision arrays, or $40n^2$ bytes. An explicit dense-inverse reference is assigned seven such arrays, or $56n^2$ bytes. In contrast, a banded Cholesky factor and a banded selected-inverse buffer together require approximately $16n(d+1)$ bytes, excluding ancillary data structures. Figure~\ref{fig:memorytradeoff} illustrates these storage relations for $d=4$.

For example, the spectral admission allowance is 640~MiB at $n=4096$. At $n=8192$ it rises to 2.5~GiB, exceeding the 2~GiB dense allocation limit imposed in the study. The latter sample length was not evaluated in the present comparison. The curves therefore represent \emph{implementation-specific memory allowances and illustrative buffer sizes}, not measurements of peak resident memory or empirical allocation failures.

\begin{figure}[t]
\centering
\includegraphics[width=\linewidth]{\VizDir/figures/fig08_memory_scaling.pdf}
\caption{Analytical source-code memory allowances as a function of series length for dense spectral SVD, an explicit inverse and two fourth-order banded working arrays. The shaded region denotes examined sample lengths and the dashed line the 2~GiB dense allocation policy. These curves are storage-policy calculations, not measured peak memory.}
\label{fig:memorytradeoff}
\end{figure}

There were 970 resource-budget exclusions in the recorded timing outcomes: 776 for the explicit-inverse reference and 194 for SVD. These were induced by prescribed computational-work thresholds, not observed out-of-memory errors. Thus the memory analysis establishes a concrete feasibility trade-off without attributing numerical failures to unsupported memory measurements.
